\section{Discussion}
\label{sec:discussion}

\subsection{Main findings}
Across the benchmark the FFNO is the strongest direct hydrostatic surrogate, reaching a dimensionless relative $L_2$ of $0.168$ against $0.235$ for the base FNO, $0.367$ for WNO, $0.475$ for the U-Net, and $0.607$ for the plain CNN (Table~\ref{tab:accuracy}). The base FNO still emulates all three numerical references to a comparable dimensionless fidelity ($0.21$--$0.26$). The speedups over the unoptimized CPU reference solvers are large, from $\sim$$8.5 \times 10 ^ {3}$ for U-FNO against the hydrostatic solver to $\sim$$1.9 \times 10 ^ {5}$ for the lightweight CNN against the same solver (Table~\ref{tab:speed}); for the FNO same-solver rows, the speedup tracks the cost of the target solver. The single-pass models' main weakness is per-frame drift over the forecast horizon (Table~\ref{tab:perframe}), which the seeded autoregressive windowed forecasters directly target: Window-FNO-5 cuts the full-trajectory error to $0.077$, and Window-FFNO-5 reaches $0.052$, at the cost of sequential inference and a true-first-frame seed that makes the comparison only partially like-for-like (Section~\ref{subsec:res-window}). Robustness is mixed but interpretable: metadata-filtered source-family diagnostics are mostly close to the in-distribution band, strict held-out-family stress tests range from mild trench degradation to a rough-source failure, and zero-shot transfer between natively simulated resolutions collapses for both the single-pass and windowed models. The deep ensemble gives a usable but overconfident uncertainty (Section~\ref{subsec:res-uncertainty}): the error--uncertainty correlation is stable around $0.65$--$0.70$ in- and out-of-distribution, while interval coverage still falls short of nominal, so the uncertainty is best read as a relative ranking signal rather than a calibrated probability.

\subsection{Effect of solver fidelity}
The two shallow-water references behave almost interchangeably from the surrogate's point of view. Their raw solver-vs-solver gap is small (relative $L_2$ $0.164$, RMSE $0.00173$; Table~\ref{tab:solvergap}), and the FNO trained on either one reaches a similar same-solver accuracy ($0.235$ hydrostatic, $0.263$ MUSCL-HR). The Boussinesq reference is the genuinely different target: its weakly dispersive correction puts it an order of magnitude away from both shallow-water solvers in physical RMSE ($\approx 0.016$--$0.017$), and that cross-family gap dwarfs the within-family one. This separation is exactly what motivates a multi-reference benchmark: a surrogate that looks solver-agnostic within the shallow-water family is in fact being tested against a near-identical target, and only the dispersive reference exposes how much of the learned behavior is tied to the training physics. The cross-solver diagnostic does \emph{not} support any emulator-superiority-like behavior: the emulator-superiority ratio exceeds one in all four cross-reference pairs (Table~\ref{tab:superiority}), from $1.76$--$1.98$ within the shallow-water family to $1.33$ across families. A surrogate reproduces the solver it was trained on far better than it reproduces a different solver, and it does not close the physical gap between solvers. We read this as the honest negative outcome of the diagnostic: the surrogate is an emulator of its training solver, not a more accurate physics model than that solver.

\subsection{Model behavior}
The accuracy ranking of the direct models is consistent with the inductive-bias argument of Section~\ref{sec:models}. The wavefields here are governed by long-range propagation: a disturbance at the source influences distant cells within the forecast horizon, so a model that can move information globally in one layer has an advantage over one that must build up an effective receptive field through depth. The FFNO and FNO spectral paths provide that global coupling, the encoder--decoder U-Net reaches intermediate range through downsampling, and the plain CNN is purely local; the measured ordering follows that reasoning. We present this as a plausible interpretation rather than a demonstrated mechanism. FFNO's direct hydrostatic improvement is not due to a larger parameter budget: it has $0.411$M parameters versus $5.443$M for the base FNO. The mode ablations also show that simply increasing retained Fourier modes from $14$ to $20$ is not enough to improve the base FNO. Where the global modeling helps least is the high-wavenumber regime: the rough source family is the hardest strict source holdout, consistent with the truncated spectral path struggling with content those sources inject. The shared failure across architectures is temporal drift in the single-pass setting and a collapse under zero-shot native-resolution transfer, neither of which is specific to the spectral model; the windowed rollout addresses the first but not the second.

\subsection{Boussinesq comparison reference}
The Boussinesq reference earns its place in the benchmark precisely because it is the one target the shallow-water surrogates cannot quietly absorb. Read in dimensionless relative $L_2$ the Boussinesq FNO looks like the best of the three ($0.210$), and its per-frame error is essentially flat over the horizon ($0.207 \to 0.208$) where the shallow-water surrogates drift upward. But this ranking is unit-dependent and we are careful not to overread it: the Boussinesq target normalization scale is about $2.3 \times$ larger, so in physical surface-elevation units the Boussinesq surrogate is actually the \emph{least} accurate (RMSE $0.00468$ vs $0.00227$ hydrostatic), and its flatness reflects the smoothness of the weakly dispersive elevation field rather than a better model. What the Boussinesq baseline adds is a controlled, weakly dispersive point of contrast that is still a synthetic NumPy reference, not real-ocean truth. It lets us measure how far a shallow-water-trained surrogate is from a different but related physics, and it shows that the cross-family gap is real and is not closed by the surrogate. We deliberately keep hydrostatic reconstruction as the primary reference and treat MUSCL-HR and Boussinesq as comparison baselines, so the Boussinesq result is framed as a fidelity-contrast diagnostic, not as a claim that the dispersive solver is closer to nature.

\subsection{Practical implications}
The benchmark is a controlled study of surrogate behavior, and we are careful not to overstate what it implies for real systems. It is not an operational tsunami forecasting tool and does not provide site-specific hazard assessment. What it can speak to is narrower: how neural surrogates of solver-generated shallow-water wavefields behave under solver mismatch, distribution shift, resolution change, a synthetic-source transfer onto fully wet real-bathymetry morphologies, and how their predictive uncertainty responds in those settings.

With that scope in mind, while the benchmark is not intended for operational forecasting, fast surrogate models of this kind may eventually contribute to scenario exploration, educational tools, rapid sensitivity analysis, or decision-support workflows, but only when combined with validated physical models and real-world data. We therefore frame any downstream relevance as a direction for future research toward such systems rather than a capability the present benchmark already provides; the 10-crop real-bathymetry morphology suite is a deliberately modest transfer diagnostic and not evidence of operational readiness.
